% Context from chapters-tex/machine-learning.tex; for mathematical reference, not compilation.
 size, this gives the empirical-risk objective~\eqref{eq-erm-param} with logistic loss:
\eql{\label{eq-logistic-optim}
	\hat\be \in \uargmin{\be \in \RR^p} 
		 E(\be)  = \frac{1}{n} \sum_{i=1}^n \loss(\dotp{x_i}{\be},y_i)  
}
where the logistic loss reads
\eql{\label{eq-logistic-loss}
	 \loss( s,y ) \eqdef \log( 1+\exp(-sy) ). 
}
Problem~\eqref{eq-logistic-optim} is a smooth convex minimization. If $X$ has full column rank, $E$ is strictly convex, so it has at most one minimizer. A finite minimizer need not exist for separable data. A positive quadratic penalty guarantees existence and uniqueness.

\begin{figure}
\centering
\includegraphics[width=.4\linewidth]{ml/classif/losses-corrected}
\caption{\label{fig-losses}
Classification losses as functions of the negative margin $s=-y\dotp{x}{\be}$, with zero margin counted as an error. The unscaled logistic loss need not upper-bound the 0--1 loss; division by $\log 2$ gives the upper-bound normalization used above.
}
\end{figure}

Figure~\ref{fig-losses} compares the 0--1, logistic, and hinge losses.

   
\paragraph{Gradient descent method.}

Write the objective as
\eq{
	 E(\be) = \Loss(X \be,y) \qwhereq \Loss(s,y)= \frac{1}{n}  \sum_i \loss(s_i,y_i), 
}
Its gradient is
\eq{
	 \nabla E(\be) = X^* \nabla \Loss(X \be,y)
      \qwhereq
      \nabla\Loss(s,y)=-\frac{y}{n}\odot\th(-y\odot s),
}
where $\odot$ is the pointwise multiplication operator, i.e. \texttt{.*} in Matlab.
 
Once $\be^{(\ell=0)} \in \RR^{p}$ is initialized (for instance at $0_{p}$), one step of gradient descent~\eqref{eq-grad-desc} reads
\eq{
	 \iit{\be} = \it{\be} - \tau_\ell \nabla E(\it{\be}). 
}

\begin{figure}
\centering
\includegraphics[width=.4\linewidth]{ml/classif/classes-2-separation-influ}
\caption{\label{fig-separation-influ}
Effect of class separation on predicted classification probabilities.
}
\end{figure}

Figure~\ref{fig-separation-influ} illustrates the method on samples from a Gaussian mixture, whose overlap is controlled by an offset $\omega$.
 
Overlaying the data on the probability map highlights the separating hyperplane $\enscond{x}{\dotp{\be}{x}=0}$.




                                                            
\subsection{Multiclass Logistic Classification}
\label{sec-multiclass-logit}

For multiclass logistic classification, introduce one weight vector for each of the
$k$ classes. Collect the vectors $\be = ( \be_\ell )_{\ell=1}^k$ as columns of $\be \in \RR^{p \times k}$.

For an input $x \in \RR^p$, convert the class scores into probabilities
using the softmax model:
\eql{\label{eq-multiclass-logitmap}
	 f(x,\be) = \pa{ \frac{ e^{\dotp{x}{\be_\ell}} }{ \sum_m e^{\dotp{x}{\be_m}} } }_\ell
}
The vector $f(x,\be)\in[0,1]^k$ gives the probabilities that $x$ belongs
to each class, and satisfies $\sum_\ell f(x,\be)_\ell=1$.

Estimate $\be$ by maximizing the log-likelihood
over the training sample:
 \eq{
	 \umax{\be \in \RR^{p \times k}} \frac{1}{n} \sum_{i=1}^n \log( f(x_i,\be)_{y_i} ) 
}
where $y_i \in \Yy = \{1,\ldots,k\}$ is the label of
observation $x_i$.

This is conveniently rewritten as
\eq{
	 \umin{\be \in \RR^{p \times k}} \Ee(\be)\eqdef\frac1n\left(\sum_i\text{LSE}(X\be)_i-\dotp{X\be}{D}\right)
}
where $D \in \{0,1\}^{n \times k}$ is the one-hot class matrix
\eq{
	  D_{i,\ell} = \choice{
          1 \qifq y_i=\ell, \\
          0 \text{otherwise}.
      }
}
and LSE is the log-sum-exp operator
\eq{
	 \text{LSE}(S)_i=\log\left(\sum_\ell\exp(S_{i,\ell})\right),\qquad\text{LSE}(S)\in\RR^n.
}
For $k=2$, identify the first class with label $+1$ and the second with label $-1$. The model then agrees with binary logistic classification in Section~\ref{sec-two-class-logit}, with parameter $\be_1-\be_2$. A single weight vector therefore suffices in that case.

Direct evaluation of log-sum-exp can
overflow when $S_{i,\ell}$ is large. A stable evaluation
subtracts the largest entry in each row before exponentiation,
using the identity
\eq{
	\text{LSE}(S+a)=\text{LSE}(S)+a
} 
for a rowwise constant $a$. This shift is especially useful for well-separated classes, whose parameter vectors $\be_\ell$ may grow large.


The gradient of the sum of the rowwise log-sum-exp values is the softmax operator
\eq{
	  \nabla_S\sum_i\text{LSE}(S)_i = \text{SM}(S) \eqdef
      \pa{
          \frac{
                  e^{S_{i,\ell}}
              }{
                  \sum_m e^{S_{i,m}}
              } }   
}
As with log-sum-exp, subtract the largest entry in each row before evaluating the exponentials.

\begin{figure}
\centering
\begin{tabular}{@{}c@{\hspace{5mm}}c@{\hspace{5mm}}c@{}}
\includegraphics[width=.25\linewidth]{ml/classif/digits}&
\includegraphics[width=.25\linewidth]{ml/classif/digits-2d}&
\includegraphics[width=.25\linewidth]{ml/classif/digits-3d}
\end{tabular}
\caption{\label{fig-digits}
Digit images and their two- and three-dimensional PCA projections.
}
\end{figure}

Once the matrix $D$ is computed, the gradient of $\Ee$ is computed as
\eq{
	 \nabla \Ee(\be) =  \frac{1}{n} X^* ( \text{SM}(X \be) - D ).  
}
The objective $\Ee$ can then be minimized by gradient descent.

To illustrate the method, we use a dataset of $n$ images of size $p = 8 \times 8$, representing digits from 0
to 9 (so there are $k=10$ classes).
 
Figure~\ref{fig-digits} displays a few representative examples as well as 2-D and 3-D PCA projections.
 
Figure~\ref{fig-digits-classes} visualizes the fitted class probabilities $f(x,\hat\be)$ by assigning colors on a regular image grid.

\begin{figure}
\centering
\includegraphics[width=.35\linewidth]{ml/classif/digits-classes}
\qquad
\includegraphics[width=.44\linewidth]{ml/classif/digits-classes-single}
\caption{\label{fig-digits-classes}
Digit classification results.
Left: probability $f(x,\hat\be)_\ell$ for each of the first nine classes, displayed in a 2-D PCA projection.
Right: colors encode the class-probability vector $f(x,\hat\be)$.
}
\end{figure}




                                                  
                                                  
                                                  
\section{Kernel Methods}
\label{sec-kernel-methods}

Linear models may fail to capture the geometry needed for a regression or classification task. Moreover, inputs such as strings and graphs need not belong to a vector space.

Kernel methods introduce a feature map into a richer space and apply linear prediction there. This produces nonlinear regression and classification rules while retaining linear-system solvers and convex optimization. The kernel trick expresses the computations through similarities between the $n$ observations, without constructing the feature vectors explicitly. Such methods are nonparametric when the number of fitted coefficients grows with $n$, allowing the model to become more expressive as data accumulate.

Algorithms that use the observations $x_i$ only through inner products can o